% !TEX root = ../main.tex
% !TEX program = lualatex
\documentclass[../main.tex]{subfiles}
\IfSubfilesClassLoaded{\externaldocument{\subfix{../build/main}}}{}
% =============================================================================
% APPENDIX: KEY ROLES IN ACQUISITION AND MODERNIZATION
% DESCRIPTION: Reference for acquisition/modernization roles and responsibilities.
% =============================================================================
\begin{document}
%====================
% Contents here
%====================

%====================
% Contents
%====================
% === ROLES ===
\IfSubfilesClassLoaded{\chapter{Key Roles and Organizations}}{\section{Key Roles in Acquisition and Modernization}}
% Program Authority (PA), Technical Authority (TA), and Contracting authority remain distinct.
\begin{description}[style=nextline, labelsep=0.5em, font=\bfseries]
  \item[\textbf{Department acquisition governance.}]
        Sources: \ac{secnavinst}~5400.15D; \ac{secnavinst}~5000.2G; DoWI~5000.85.
        \begin{itemize}
          \item \ac{asnrdanda}, the \ac{sae} and \ac{cae} for the \ac{don}, sets acquisition policy, charters \acp{peo} / \acp{drpm}, and serves as or delegates \ac{mdauthority} for \ac{don} programs according to category and designation.
          \item \acp{peo} and \acp{drpm} hold delegated program authority, establish baselines, and are accountable for translating capability needs into executable acquisition strategies.
          \item Current as of 2026-06-04, \ac{don} has publicly established nine \acp{pae}. PAE Mission Systems spans PEO C4I, PEO Digital, PEO IWS, PEO MLB, DRPM Overmatch, DRPM LRNFO, DRPM MILE, Minotaur from NAVAIR's PMA-290 program office, and mission-system elements across NAVWAR, NAVSEA, NAVAIR, and Marine Corps Systems Command~\autocite{DON-PAE-Expansion-May2026,DON-PAE-Mission-Systems-2026}.
          \item \ac{syscom} Commanders provide the workforce, infrastructure, and \ac{ta} warrants that underpin \ac{peo} execution inside the Navy matrix construct.
        \end{itemize}

  \item[\textbf{Joint requirements governance.}]
        Sources~\autocite{CJCSM5123-01,jcids-manual-2021,secwar-acq-transformation-2025}.
        \begin{description}
          \item \textbf{\ac{jroc}:} Under the current JFRP, remains the senior joint requirements body for Joint Force Design, Joint Capability Integration, Combatant Command requirements, \ac{cpm}, and by-exception reviews; legacy \ac{jrocm} direction remains binding unless formally relieved.
          \item \textbf{Joint Staff J-8/\ac{jcb}:} Supports \ac{jroc} deliberations, JFRP staffing, and joint-impact integration; do not answer that most Navy Service/component requirements wait for default \ac{jcb} validation before acquisition execution.
          \item \textbf{\acp{fcb}:} Provide functional analysis and integration support for JFRP issues, performance-attribute certifications or endorsements, and cross-functional joint impacts rather than serving as a universal Service-document validation layer.
          \item \textbf{\acp{pae} and \acp{ctctrade}:} The 2025 Acquisition Transformation Strategy empowers \acp{pae} to convene Capability Trade Councils that record requirement/cost/schedule trades in real time so \ac{peo} teams keep the \ac{das} synchronized with evolving demand signals.
        \end{description}

  \item[\textbf{Combatant command sponsors.}]
        Sources~\autocite{edo-2-1-1-command-structures-2025,jcids-manual-2021}.
        \begin{itemize}
          \item \textbf{\ac{ccdr}s} own warfighting problems, sign \acp{juon}/\acp{jeon}, and direct lead Services or agencies to prosecute capability gaps that demand rapid solutions.
          \item \textbf{\ac{ccmd} staffs} conduct mission-thread analysis and architecture sprints with Fleets, labs, and industry to translate operational risk into measurable capability requirements before Component approval, Joint Staff integration review, or urgent/emergent validation as applicable.
          \item Commanders remain voting participants during validation reviews and oversee fielding to ensure delivered capability solves the theater problem that generated the requirement.
        \end{itemize}

  \item[\textbf{Joint Chiefs leadership.}]
        Sources~\autocite{edo-2-1-1-command-structures-2025,jcids-manual-2021,USC-10-181}.
        \begin{description}
          \item \textbf{\ac{cjcs}:} Principal military adviser to the President/\ac{secdef}, co-chairs the \ac{jroc}, and leads joint concept and readiness development that frames every Capability-Based Assessment.
          \item \textbf{\ac{vcjcs}:} Chairs the \ac{jcb}, co-leads the \ac{dmag}, and arbitrates cross-Service requirement or resourcing disputes before they escalate to the Deputies or Secretary level.
        \end{description}

  \item[\textbf{OPNAV corporate leadership.}]
        Source~\autocite{edo-2-1-1-command-structures-2025}.
        \begin{description}
          \item \textbf{\ac{cno}:} Serves as Service Chief, chairs the Navy corporate board, and approves Navy positions for Gate~0/1 decisions before submissions move to the Joint Staff.
          \item \textbf{\ac{opnav} resource integrator, budget lane, and sponsors (\ac{n8}, \ac{n82}/\ac{fmb}, \ac{n9}, \ac{n9i}/9I):} Translate Fleet priorities into the \ac{pom}/\ac{fydp}, enforce two-pass/seven-gate rigor, and synchronize requirements, budgets, and \ac{peo} execution.
        \end{description}

  \item[\textbf{Acquisition decision authorities.}]
        Sources~\autocite{DoDD5000-01,DoDI5000-85,SECNAV5000-2G}.
        \begin{description}
          \item \textbf{\ac{dae} / \ac{usdas}:} Issues acquisition policy, chairs the Defense Acquisition Board, and retains \ac{mdauthority} for \ac{acat}~ID/\ac{iam} programs unless delegated.
          \item \textbf{\ac{sae}/\ac{asnrdanda}:} Assigns Navy \ac{mdauthority} for \ac{acat}~II and below, approves Service-tailored governance (e.g., two-pass/seven-gate), and charters \acp{peo}/\acp{drpm}.
          \item \textbf{Delegated \ac{mdauthority} (\ac{peo}/\ac{drpm}):} Signs Acquisition Decision Memoranda, manages Configuration Steering Boards, and keeps \ac{kpp}/\ac{apb} trades synchronized with requirements direction and statutory certifications.
        \end{description}

  \item[\textbf{Program Authority chain.}]
        Sources: \ac{secnavinst}~5000.2G; DoWI~5000.85.
        \begin{description}
          \item \textbf{\ac{peo}/\ac{pa}}: Owns program outcomes, charters \acp{pm}, approves acquisition strategies and baselines.
          \item \textbf{\ac{drpm}}: Direct-report program leads for special access or priority portfolios with the same milestone authority as \acp{peo}.
          \item \textbf{\ac{spm}/\ac{pm}}: Accountable for cost/schedule/performance, leads risk management, orchestrates \acp{ipt}.
          \item \textbf{\ac{shapm}}: Delivers new hulls; coordinates design/build/activation sequencing.
          \item \textbf{\ac{slm}}: Drives sustainment and modernization packages for in-service assets.
          \item \textbf{\ac{parm}/\ac{spd}}: Sponsors platform-level upgrades and alteration packages.
        \end{description}

  \item[\textbf{Technical Authority chain.}]
        Source: \ac{secnavinst}~5400.15D.
        \begin{description}
          \item \textbf{\ac{uta} / \ac{cheng}}: Issues enterprise technical policy, owns warrants.
          \item \textbf{\acp{twh}}: Certifies design compliance, adjudicates departures/concurrence packages.
          \item \textbf{Engineering Agents (\ac{isea}, \ac{da}, \ac{aea}, \ac{sia}, \ac{tda})}: Execute lifecycle engineering; provide in-service engineering decisions under delegated authority.
        \end{description}

  \item[\textbf{Operational test oversight.}]
        Sources~\autocite{DoDI5000-85,DoDAcqGuidebook}.
        \begin{description}
          \item \textbf{\ac{dote}:} Independent \ac{osd} authority that concurs on \ac{temp}s, monitors operational test realism, and reports suitability/effectiveness outcomes to Congress.
          \item[\textbf{Service Operational Test Agencies.}] \ac{comoptevfor} (\acs{cotf}), \ac{afotec}, \ac{atec}, and \ac{mcotea} plan and execute \ac{iote}/\ac{fote}, render the \ac{otrr} recommendation, and provide Service suitability findings to the \ac{mdauthority}.
        \end{description}
  \item[\textbf{Operational test readiness and fleet release.}]
        Sources~\autocite{edo-3-5-6-test-eval-part-ii-2025,DoDI5000-89}.
        \begin{description}
          \item \textbf{\ac{cotf}/\ac{comoptevfor}:} Chairs \ac{otrr} events, certifies whether \ac{ote} adequacy and logistics/training enablers support Fleet introduction, and can recommend restricted Fleet release when Cat~I deficiencies remain.
          \item \textbf{\ac{mdauthority}/\ac{peo}:} Disposition Cat~I/II findings, resource fixes, and may defer \ac{frpdr} until \ac{ote} adequacy and deficiency burn-downs meet Title~10 standards.
          \item \textbf{\ac{tycom}:} Accepts Fleet risk, sets operating restrictions, and updates tactics and training based on \ac{ote} and \ac{lfte} findings.
        \end{description}
  \item[\textbf{Developmental test governance.}]
        Sources~\autocite{DoDI5000-89}.
        \begin{description}
          \item \textbf{\ac{dasdte}:} Designates Chief Developmental Testers and Lead Developmental Test and Evaluation Organizations for \ac{acat}~I programs, adjudicates test resource shortfalls, and co-chairs the Defense T\&E \ac{wipt}.
          \item \textbf{\ac{cdt}:} Synchronizes the integrated test schedule, deficiency management cadence, and readiness evidence required for each \ac{trr}/milestone.
          \item \textbf{\ac{ldto}:} Provides instrumentation, modeling \& simulation accreditation, range execution, and authoritative data packages for developmental tests.
        \end{description}
  \item[\textbf{Live-fire statutory leads.}]
        Sources~\autocite{USC-10-4172,DoDI5000-89}.
        \begin{description}
          \item \textbf{\ac{dote} Live Fire Test Director:} Approves plans, monitors execution realism, and delivers the statutory lethality/vulnerability report to Congress.
          \item \textbf{Program Manager/\ac{cheng}:} Funds threat-representative articles, integrates survivability design fixes, and aligns \ac{lfte} data with the Acquisition Program Baseline and sustainment plans.
          \item[\textbf{Service Safety Centers and Warfare Centers:}] Provide casualty data, threat surrogates, and post-test forensic analysis that feed both live-fire reports and Fleet tactics updates.
        \end{description}

  \item[\textbf{NAVSEA headquarters roles.}]
        Source: \ac{edo} Coursebook Module~2.1.2 (\ac{navsea} Organization, 2025 edition).
        \begin{description}
          \item \textbf{\ac{sea}~01 Comptroller}: \ac{bso} lead for \ac{navsea} financial governance and funds control.
          \item \textbf{\ac{sea}~02 Contracts}: Enterprise contracting authority, policy, and warrant management.
          \item \textbf{\ac{sea}~03 Cyber Engineering and Digital}: Drives cyber resiliency, digital engineering, and data transformation initiatives.
          \item \textbf{\ac{sea}~04 Industrial Operations}: Oversees public shipyards, maintenance execution, and quality assurance.
          \item \textbf{\ac{sea}~05 \ac{cheng}}: Chief \ac{ta}; issues warrants, certifies designs, and maintains technical standards.
          \item \textbf{\ac{sea}~06 Sustainment}: Manages product support strategies and lifecycle logistics integration.
          \item \textbf{\ac{sea}~07 Undersea Warfare}: Dual-hatted as \ac{peouws}; oversees undersea combat systems sustainment.
          \item \textbf{\ac{sea}~08 Nuclear Propulsion}: Three-hatted nuclear propulsion authority across \ac{don} and \ac{doe} roles.
          \item \textbf{\ac{sea}~09 Safety and Regulatory Compliance}: Aligns enterprise safety governance and reporting.
          \item \textbf{\ac{sea}~10 Total Force and Corporate Ops}: Manages workforce planning, corporate services, and governance.
          \item \textbf{\ac{sea}~21 In-Service Ships/\ac{cnrmc}}: Leads surface-ship sustainment and modernization (\ac{pms}~321/326/339/421/443/451, \ac{sea}~21I, \ac{surfmepp}).
        \end{description}

  \item[\textbf{NAVWAR enterprise directorates.}]
        Source: \ac{edo} Coursebook Module~2.1.3 (\ac{navwar} Enterprise, 2024 edition).
        \begin{description}
          \item[\textbf{Code~1.0 Comptroller}] : Budget formulation/execution, \ac{bso} duties, and funds certification.
          \item[\textbf{2.0 Contracts}] : Contracting policy, strategy reviews, and award/administration oversight.
          \item[\textbf{3.0 Counsel}] : Acquisition law, ethics, protests, and claim resolution.
          \item[\textbf{4.0 Logistics and Fleet Support}] : Lifecycle logistics, technical data, and fleet distance support integration.
          \item \textbf{5.0 \ac{cheng}/\ac{ta}}: Enterprise architectures, interoperability, and cyber certification authority.
          \item[\textbf{6.0 Program Management}] : Portfolio governance, milestone preparation, and \ac{peo} integration.
          \item[\textbf{7.0 Science and Technology}] : S\&T portfolio management, prototyping, and technology transition.
          \item[\textbf{8.0 Corporate Operations}] : Workforce, \ac{cio} services, facilities, security, and public affairs.
          \item \textbf{\ac{frd}-100 Fleet Support}: Deployed engineering assistance and sustainment response.
          \item \textbf{\ac{frd}-200 Installations}: Shore/afloat \ac{c4i} installation planning and cutover execution.
        \end{description}

  \item[\textbf{Program Executive Offices (Navy portfolios).}]
        Sources: \ac{secnavinst}~5400.15D; \ac{edo} Coursebook Modules~2.1.4 and~3.1.5.
        \begin{description}
          \item \textbf{\ac{peocvn}}: Designs, builds, and sustains nuclear-powered aircraft carriers (e.g., \ac{pms}~312/378/379).
          \item \textbf{\ac{peoiws}}: Develops and sustains ship, submarine, aircraft, and allied combat-system portfolios; mission-aligned \ac{iws} directorates cover sensors, weapons, C2, and allied integration.
          \item \textbf{\ac{peoships}}: Oversees surface combatant and amphibious ship construction and modernization.
          \item \textbf{\ac{peousc}}: Leads \ac{lcs}, \ac{ffg}~62, expeditionary, and unmanned surface/undersea portfolios.
          \item \textbf{Team Submarines (\ac{peossn}, \ac{peossbn}, \ac{peocolumbia}, \ac{peouws})}: Manages attack/strategic submarine acquisition, in-service support, and undersea combat systems.
          \item \textbf{\ac{peoc4i}}: Delivers fleet \ac{c4i}; \ac{pmw}~1XX focus on capability development, \ac{pmw}~7XX on platform integration.
          \item \textbf{\ac{peodigital}}: Provides enterprise digital services (e.g., Flank Speed) across the \ac{don}.
          \item \textbf{\ac{peomlb}}: Modernizes manpower, logistics, and business IT systems.
          \item \textbf{PAE Mission Systems cue}: Treat PAE Mission Systems as current acquisition-governance alignment across NAVWAR, NAVSEA, NAVAIR, and Marine Corps Systems Command mission-system portfolios; retain current public PEO, PMW, and IWS labels for board recall until official source pages replace them.
        \end{description}

  \item[\textbf{Contracting authority (KO family).}]
        Sources: \ac{far}; \ac{dfars}; \ac{nmcars}; \ac{edo} Coursebook Module~3.2.1.
        \begin{description}
          \item \textbf{\ac{pco}}: Plans the acquisition, synchs with \acp{pm} before solicitation, awards and signs contracts/mods.
          \item \textbf{\ac{aco}}: Oversees post-award performance, surveillance, and payment/contract administration (often \ac{dcma} / \ac{navsea} field activities).
          \item \textbf{\ac{tco}}: Leads partial/full terminations, settlement negotiations, and equitable adjustments.
          \item \textbf{\ac{ko} warrant}: Defines the dollar/authority limits; only the warranted \ac{ko} can bind or obligate the Government.
          \item \textbf{\ac{cor} / assistant \ac{pm} / engineering support}: Provide technical surveillance and acceptance recommendations; cannot direct work or obligate funds (\ac{far} Parts~1 and~42).
          \item \textbf{\ac{hca}}: Approves actions such as letter contracts and high-value single-award \acp{idiq}; may delegate no lower than flag/Senior Executive levels within the \ac{don} per \ac{far}~1.601 and \ac{dfars}/\ac{nmcars} supplements.
          \item \textbf{\ac{sae}, \ac{asnrdanda}}: Signs \acp{df} for multiyear contracting, extraordinary relief, and other actions reserved to the Service Acquisition Executive under \ac{far} Subparts~1.7 and~17.1.
        \end{description}

  \item[\textbf{Program Manager / KO partnership.}]
        Sources: DoWI~5000.85; \ac{edo} Coursebook Module~3.2.1.
        \begin{itemize}
          \item \ac{pm} integrates warfighter need, technical baseline, and budget; \ac{ko} \ac{cica}, incentives, and change management before \ac{rfp} release or modification execution.
        \end{itemize}

  \item[\textbf{Source selection governance.}]
        Sources: \ac{far} Parts~1, 5, and~15; \ac{dod} Source Selection Procedures; \ac{navsea} Source Selection Guide (2022); \ac{edo} Coursebook Module~3.2.2.
        \begin{description}
          \item \textbf{\ac{ssa}}: Senior official who approves the Source Selection Plan, receives \ac{ssac}/\ac{sseb} recommendations, and signs the best-value decision memorandum.
          \item \textbf{\ac{ssac}}: Advisory council, when established or required by \ac{dod}/agency source-selection procedure, that synthesizes \ac{sseb} findings, compares proposals across factors, and briefs the \ac{ssa} on trade-offs.
          \item \textbf{\ac{sseb}}: Multi-disciplinary evaluators (technical, management, past performance, cost/price) who rate proposals against the solicitation evaluation factors, including Section~M when the uniform contract format is used.
          \item[\textbf{Small Business Professional / Competition Advocate}] : Confirms set-aside decisions, reviews subcontracting plans, and endorses synopsis waivers.
          \item[\textbf{Cost/Price Analyst \& Legal Counsel}] : Validate reasonableness determinations, alignment of Sections~L/M, and clause sufficiency before release.
          \item[\textbf{Award Fee Determining Official \& Award Fee Board}] : Plans award-fee periods, chairs performance reviews, and signs the determination memo authorizing or withholding fee payments (\ac{far} Part~16; \ac{navsea} Source Selection Guide).
        \end{description}

  \item[\textbf{Fiscal control \& certification.}]
        Source: \ac{dodfmr}~7000.14-R Volume~3.
        \begin{description}
          \item[\textbf{Funds Certifying Official/Comptroller}] : Verifies purpose/time/amount before obligation; \ac{ada} safeguard.
          \item[\textbf{Program/Budget Analyst}] : Tracks execution, monitors reprogramming thresholds, prepares obligation/expenditure burn-down.
          \item[\textbf{Resource allocation chain}] : \ac{ousdc} apportions to the Navy; \ac{opnav} (\ac{n82}/\ac{fmb}) allocates to \acp{bso}; \ac{syscom} comptrollers (e.g., \ac{sea}~01) issue suballocations to executing activities.
        \end{description}

  \item[\textbf{PPBE resource sponsors and integrators.}]
        Sources: \ac{edo} Coursebook Module~3.1.1 (\ac{ppbe}, 2025 edition), NAVMAC Activity Manpower Management Guide Section~24, and current public \ac{opnavinst} readiness/assessment instructions.
        \begin{description}
          \item \textbf{\ac{cape}}: Independent analytic staff supporting Program Review decisions and \ac{ice}/\ac{poe} for major programs~\autocite{USC-139c-CAPE,edo-3-1-1-ppbe-2025}.
          \item \textbf{\ac{n8}}: Navy PPBE resource integrator (``checkbook''): builds a balanced Navy \ac{pom} across portfolios; do not confuse this with the NAVMAC resource-sponsor code list.
          \item \textbf{\ac{n80}}: Runs the \ac{pom} build and integration; adjudicates sponsor submissions.
          \item \textbf{\ac{n81}}: Provides force-structure and analytic assessments supporting \ac{pom} trades.
          \item \textbf{\ac{n82}}: Budget formulation/execution chain (\ac*{fmb}); manages allocation, execution controls, and reprogramming posture.
          \item \textbf{\ac{n83}}: Fleet Readiness; \ac{opnav} point of contact for ship maintenance and fleet-readiness matters, and lead on CNO N8 ship-maintenance responsibilities.
          % \item \textbf{N84/N89}: Treat as legacy or internally verified labels only; the current public source set used here does not support presenting them as current N8 resource offices.
          \item \textbf{N7}: Warfighting Development in the NAVMAC resource-sponsor code list; do not confuse this with the CNO principal-assistant N3/N5 operations, plans, and strategy lane.
          \item \textbf{\ac{n9}}: Warfare-systems sponsor family (``demand signal''): validates and advocates warfare requirements.
          \item \textbf{\ac{n9i}/9I}: Cross-domain warfare integration; current public manpower coding lists this function as 9I Warfare Integration. Treat N91 as an older or alternate public-reference label, not the current resource-sponsor code.
          \item \textbf{N94}: Test and Evaluation sponsor in the public NAVMAC resource-sponsor code list.
          \item \textbf{\ac{n95}/\ac{n96}/\ac{n97}/\ac{n98}}: Expeditionary/surface/undersea/air warfare sponsors under \ac{n9}.
          \item \textbf{\ac{n99}}: Unmanned warfare systems sponsor; integrates unmanned air/surface/undersea capability portfolios and defends their investments in \ac{ppbe}.
          \item \textbf{\ac{r3b}/\ac{ncb}}: Navy governance boards that align requirements and resourcing decisions feeding the \ac{pom}.
        \end{description}

  \item[\textbf{Congressional resource chain.}]
        Source: \ac{edo} Coursebook Module~3.1.2 (Congressional Enactment, 2025 edition).
        \begin{description}
          \item \textbf{\ac{hasc} / \ac{sasc}}: Authorize defense programs and policy in the annual \ac{ndaa}.
          \item \textbf{\ac{hac} / \ac{sac}}: Produce appropriations bills that provide \ac{ba} to execute Navy programs.
          \item \textbf{\ac{hbc} / \ac{sbc}}: Set topline guidance and 302 allocations through the budget resolution process.
          \item \textbf{\ac{cbo}, \ac{crs}, \ac{gao}}: Supply independent scoring, research, and oversight that shape committee deliberations.
        \end{description}

  \item[\textbf{Cyber/Authorizations (when IT/C2 in scope).}]
        Source: \ac{secnavinst}~5000.2G.
        \begin{description}
          \item \textbf{\ac{ao}}: Grants Authority to Operate; enforces \ac{rmf} controls that drive design and integration requirements.
        \end{description}
\end{description}
%====================
% End of file
%====================
\IfSubfilesClassLoaded{
  \RenewDocumentCommand{\entryname}{}{\textbf{\color{Modern} Acronym}}
  \RenewDocumentCommand{\descriptionname}{}{\textbf{\color{Modern} Definition}}
  \printnoidxglossary[
    type=\acronymtype,
    title=Acronyms,
    style=long-booktabs
  ]}{}
\end{document}
